% BHU PYQ -- Manually transcribed & error-corrected
\documentclass[11pt,a4paper]{article}

\usepackage[a4paper,top=2.5cm,bottom=2.5cm,left=2.8cm,right=2.8cm,headheight=14pt]{geometry}
\usepackage{amsmath,amssymb}
\usepackage[T1]{fontenc}
\usepackage[utf8]{inputenc}
\usepackage{lmodern}
\usepackage{microtype}
\usepackage{fancyhdr}
\usepackage{enumitem}
\usepackage{hyperref}
\usepackage{lastpage}
\usepackage{parskip}

\hypersetup{colorlinks=true,urlcolor=black,linkcolor=black,
  pdftitle={BPT-502: Classical Mechanics -- BHU PYQ},pdfauthor={Banaras Hindu University}}

\pagestyle{fancy}\fancyhf{}
\renewcommand{\headrulewidth}{0.4pt}\renewcommand{\footrulewidth}{0.4pt}
\fancyhead[L]{\small   \textit{Banaras Hindu University}}
\fancyhead[R]{\small   \textit{BPT-502: Classical Mechanics}}
\fancyfoot[C]{\small Page  \thepage\ of \pageref{LastPage}}


\newlist{parts}{enumerate}{2}
\setlist[parts,1]{label=(\alph*),leftmargin=*,itemsep=5pt,topsep=3pt}
\setlist[parts,2]{label=(\roman*),leftmargin=*,itemsep=3pt,topsep=2pt}

\newcommand{\pts}[1]{\hfill   \textit{[#1]}}


\begin{document}


\begin{center}
  {\large\bfseries BANARAS HINDU UNIVERSITY}\\[2pt]
  {\normalsize B.Sc. (Hons.) Semester V Examination, 2016-17}\\[6pt]
  \rule{0.8\linewidth}{0.5pt}\\[6pt]
  {\large\bfseries Paper No. BPT-502: Classical Mechanics}\\[4pt]
  {\normalsize Time: 3 Hours\hfill Maximum Marks: 70}
\end{center}
\rule{\linewidth}{0.4pt}
\smallskip



\noindent   \textit{Instructions: Answer all questions. Symbols have their usual meanings.}
\bigskip




\noindent   \textbf{Question 1.} Answer any   \textbf{four} of the following: \pts{4 $ \times$ 5 = 20}

\begin{parts}
  \item How many generalized/quasi-generalized coordinates are required to describe the following systems?
    
\begin{parts}
      \item A rigid body in 3 dimensions.
      \item A solid sphere rolling without slipping on an inclined plane.
      \item A spherical pendulum in 3 dimensions.
    \end{parts}
  \item What is isotropy of space? Prove that the conservation law of angular momentum follows from it.
  \item A particle moving in a central force field located at $r = 0$ describes a spiral orbit $r = a \theta$. Find the force law.
  \item The Lagrangian of a system is given by:
    \[ L = \frac{1}{2} m \dot{\vec{r}}^2 + \dot{\vec{r}} \cdot (\vec{a}  \times \vec{r}) + V(\vec{r}) \]
    where $\vec{a}$ is an arbitrary vector depending on time only. Find the equations of motion for the system.
  \item A particle is constrained to move on the surface of a cylinder $x^2 + y^2 = R^2$ and is subjected to a force directed towards the origin and proportional to the distance of the particle from the origin. Find the equations of motion for the particle using the Lagrange multiplier method.
\end{parts}

\medskip\rule{\linewidth}{0.2pt}




\noindent   \textbf{Question 2.}

\begin{parts}
  \item Consider a particle of mass $m$ free to slide on a smooth helical wire whose equations in cylindrical coordinates are given by $\rho = R$, $z = k \theta$.
    
\begin{parts}
      \item Construct the Lagrangian of the system. \pts{4}
      \item Find the equation of motion and solve it to discuss the motion of the particle (if it is released at rest at $\rho=R, z=z_0,  \theta= heta_0$). \pts{5}
      \item Construct the Hamiltonian of the system and find the canonical equations of motion. \pts{5}
    \end{parts}
\end{parts}
\hfill   \textbf{OR}

\begin{parts}
  \item Show that the Lorentz force on a charged particle in an electromagnetic field can be obtained from a velocity-dependent potential $U = q\phi - \frac{q}{c}\vec{v}\cdot\vec{A}$. \pts{5}
  \item Write the Lagrangian of a charged particle in the electromagnetic field using this potential. \pts{4}
  \item Write down the Hamiltonian for a system consisting of an electron, a proton, and a neutron. How will the Hamiltonian be modified in the presence of a uniform magnetic field $\vec{B} = B_0 (\hat{i} + 2\hat{j})$? \pts{5}
\end{parts}

\medskip\rule{\linewidth}{0.2pt}




\noindent   \textbf{Question 3.}

\begin{parts}
  \item A particle of mass $m$ is moving on a plane under an attractive central force given by:
    \[ F(r) = -\frac{\alpha}{r^2} - \frac{\beta}{r^3}, \quad \alpha,\beta > 0 \]
    
\begin{parts}
      \item Find the effective potential for this system. \pts{3}
      \item Discuss the motion of the particle qualitatively for various total energy values of the particle. \pts{4}
      \item Show that the angular momentum of the particle is conserved. \pts{3}
      \item Show that the position vector of the particle sweeps out equal areas in equal times. \pts{4}
    \end{parts}
\end{parts}

\medskip\rule{\linewidth}{0.2pt}




\noindent   \textbf{Question 4.}

\begin{parts}
  \item Define canonical transformation. How is it useful? \pts{4}
  \item For what values of $\alpha$ and $\beta$ is the transformation:
    \[ Q = q^{\alpha} \cos(\beta p), \quad P = q^{\alpha} \sin(\beta p) \]
    canonical? \pts{4}
  \item Find a suitable generating function for this canonical transformation. \pts{6}
\end{parts}
\hfill   \textbf{OR}

\begin{parts}
  \item Express the Hamiltonian $H = \frac{p^2}{2m} + \frac{1}{2}m\omega^2 q^2$ in terms of $Q, P$ by using a canonical transformation generated by $F_1(q, Q) = \frac{1}{2} m\omega q^2 \cot Q$. \pts{7}
  \item Calculate the Poisson brackets:
    \[ \{L_x, x\}, \quad \{L_x, p_y\}, \quad  \text{and} \quad \{(\vec{r}  \times \vec{p})_z, z\} \] \pts{7}
\end{parts}

\medskip\rule{\linewidth}{0.2pt}




\noindent   \textbf{Question 5.}

\begin{parts}
  \item Show that if $A$ and $B$ are two constants of motion, then their Poisson bracket $\{A, B\}$ is also a constant of motion. \pts{6}
  \item State and prove the virial theorem. \pts{8}
\end{parts}
\hfill   \textbf{OR}

\begin{parts}
  \item A particle of mass $m$ is falling freely from a height $h$. Find the solution for the particle using the Hamilton-Jacobi equation. \pts{7}
  \item Discuss the idea of mechanical similarity. Use this to show that the time of fall for a freely falling body is proportional to the square root of its initial height. \pts{7}
\end{parts}

\vfill
\rule{\linewidth}{0.4pt}

\begin{center}\small
  \href{https://bhu-exam.vercel.app/}{Click here to visit our website}\\[4pt]
  \href{https://me-aryan.vercel.app/}{Click here to visit our contributor}\\[4pt]
  \href{https://quashan-me.vercel.app/}{Click here to visit our contributor}
\end{center}

\end{document}
